\documentclass[11pt,a4paper]{article}

% ── Packages ──────────────────────────────────────────────────────────────────
\usepackage[utf8]{inputenc}
\usepackage[T1]{fontenc}
\usepackage{amsmath,amssymb,amsthm}
\usepackage{booktabs}
\usepackage{multirow}
\usepackage{array}
\usepackage{hyperref}
\usepackage{url}
\usepackage{geometry}
\usepackage{graphicx}
\usepackage{listings}
\usepackage{xcolor}
\usepackage{enumitem}
\usepackage{authblk}
\usepackage{abstract}

\geometry{
  a4paper,
  top=25mm, bottom=25mm,
  left=25mm, right=25mm,
}

% ── Hyperref setup ────────────────────────────────────────────────────────────
\hypersetup{
  colorlinks=true,
  linkcolor=blue!70!black,
  citecolor=blue!70!black,
  urlcolor=blue!70!black,
  pdftitle={Narrative Quantification: Compiling Narrative Structures into Computable Representations for Large Cognitive Models},
  pdfauthor={Tomohiko Nakamura},
}

% ── Code listing style ────────────────────────────────────────────────────────
\lstdefinestyle{json}{
  basicstyle=\ttfamily\small,
  backgroundcolor=\color{gray!8},
  frame=single,
  rulecolor=\color{gray!40},
  breaklines=true,
  showstringspaces=false,
  keywordstyle=\color{blue!70!black},
  stringstyle=\color{orange!70!black},
  commentstyle=\color{green!50!black},
}

% ── Theorem environments ──────────────────────────────────────────────────────
\newtheorem{definition}{Definition}
\newtheorem{proposition}{Proposition}

% ─────────────────────────────────────────────────────────────────────────────
\begin{document}

% ── Title ─────────────────────────────────────────────────────────────────────
\title{\textbf{Narrative Quantification:}\\
Compiling Narrative Structures into Computable\\
Representations for Large Cognitive Models}

\author[1]{Tomohiko Nakamura}
\affil[1]{Independent Researcher, Japan \\ \texttt{tomohiko@gemminai.com}}

\date{March 2026}

\maketitle

% ── Abstract ──────────────────────────────────────────────────────────────────
\begin{abstract}
Human societies interpret events primarily through narrative structures
rather than isolated data points. This paper proposes \textbf{Narrative
Quantification}, a framework for transforming narrative information
into structured cognitive data. By decomposing narratives into actors,
events, conflicts, resolutions, emotional context, and causal
relationships, narratives can be stored as structured datasets and
analyzed computationally.

To operationalize this transformation, the paper introduces the concept
of a \textbf{Narrative Compiler}, a system that converts narrative text
into structured cognitive representations such as the \textbf{24TAG
structure}. These representations enable the creation of narrative
databases that may serve as the foundation for \textbf{Large Cognitive
Models (LCM)}---systems designed to reason over structured
representations of human narratives rather than purely probabilistic
language tokens.

This work suggests that narrative quantification may function as a new
layer of cognitive infrastructure for future AI systems.
\end{abstract}

\noindent\textbf{Keywords:} Narrative Quantification, Large Cognitive Models,
24TAG, Narrative Compiler, Cognitive Infrastructure, State Hash, JSON
Canonicalization Scheme

\hrule
\vspace{1em}

% ── 1. Introduction ───────────────────────────────────────────────────────────
\section{Introduction}

Most contemporary AI systems rely on probabilistic language modeling.
Large Language Models (LLMs) generate text by predicting token
probabilities within massive corpora. While effective for language
generation, these systems do not explicitly represent narrative
structure, causal relationships, or long-term narrative dynamics.

Human cognition, however, rarely processes information as isolated
tokens. Instead, humans organize knowledge through \textbf{narratives}:
structured relationships between actors, events, motivations, and
outcomes.

This paper proposes that quantifying narrative structures can provide an
alternative computational representation of human cognition. By
decomposing narratives into structured components, it becomes possible
to construct databases of collective narrative patterns.

% ── 2. Narrative Cognition ────────────────────────────────────────────────────
\section{Narrative Cognition}

Human cognition frequently follows a narrative processing loop:

\begin{equation}
\text{Event} \rightarrow \text{Meaning} \rightarrow \text{Narrative} \rightarrow \text{Memory}
\end{equation}

Experiences are rarely stored as raw data. Instead, individuals
interpret events through causal explanations and emotional context,
forming narratives that are later stored in memory. This mechanism
allows humans to compress complex experiences into structured stories
that \textbf{predict future outcomes and guide subsequent decisions}.
Narrative cognition therefore represents a fundamental cognitive process
rather than a mere cultural artifact.

% ── 3. Narrative Quantification ───────────────────────────────────────────────
\section{Narrative Quantification}

Narrative Quantification is the process of converting narrative
structures into analyzable data. Narratives can be decomposed into
structural components such as:

\begin{itemize}[noitemsep]
  \item Actors
  \item Events
  \item Conflict
  \item Resolution
  \item Emotional context
  \item Causal relationships
\end{itemize}

Once extracted, these components can be represented through structured
tags and stored within databases. This enables large-scale analysis of
narrative dynamics across media, politics, economics, and culture.

\subsection{Decomposition}

The first stage of quantification is \textbf{Decomposition}, where raw
narrative text is segmented into its primary semantic units. Rather than
treating a story as a continuous stream of words, we identify discrete
entities and their roles. By isolating \textbf{Actors} (the agents of
change), \textbf{Events} (the actions taken), and \textbf{Causal
Relationships} (the `why' behind the actions), we transform linguistic
ambiguity into a set of distinct variables.

\subsection{Event Formation}

The second stage of quantification is \textbf{Event Formation}, where
continuous occurrences are discretized into logical structures. Within
this framework, \textbf{an event is defined as the minimum causal
interaction unit between actors unfolding over time.}

Unlike traditional data processing, event formation recognizes the
inherent unity of an action and its subsequent response. A typical
event unit is structured as follows:

\begin{itemize}[noitemsep]
  \item \textbf{Actors:} The agents initiating or receiving the interaction.
  \item \textbf{Action-Reaction Pair:} The core causal loop.
  \item \textbf{Temporal Boundary:} The specific window from trigger to outcome.
\end{itemize}

By defining the event as the \emph{minimum causal interaction unit}, we
enable LCMs to perform reasoning not just on what happened, but on the
\textbf{logic of the interaction itself}.

\subsubsection{Eventualization: Entropy Suppression through Gravitational Convergence}

Without event formation, information entropy increases monotonically over
time. A single real-world occurrence typically generates dozens of
articles across different media outlets — each carrying its own framing,
emphasis, and omissions. Left unstructured, this proliferation of
fragments causes the informational representation of a single event to
\emph{diverge}, consuming disproportionate computational and cognitive
resources.

\textbf{Eventualization} is the process by which the \textbf{Event
Builder} counteracts this divergence. Rather than treating each article
as an independent data point, the Event Builder applies a gravitational
model: semantically related articles are attracted toward a shared
\texttt{event\_id}, collapsing the high-entropy cloud of fragments into
a single, stable coordinate in Narrative Space.

Formally, given a set of $n$ articles $\{a_1, a_2, \ldots, a_n\}$
referencing the same real-world occurrence, Eventualization assigns them
to a common event object $E$:

\begin{equation}
  E = \text{EventBuilder}\!\left(\{a_1, \ldots, a_n\}\right), \quad
  \text{sim}(e(a_i),\, e(a_j)) \geq \theta
\end{equation}

where $e(\cdot)$ is the embedding function mapping an article to its
vector representation in Narrative Space, and $\theta$ is the similarity
threshold above which two articles are considered to reference the same
event. This operation is the narrative analogue of entropy reduction:
the $n$ fragments, each carrying partial and often contradictory
information, are compressed into one canonical object whose
informational content is greater than the sum of its parts.

\subsubsection{The Centroid Vector and Crystallization Completeness}

A direct consequence of Eventualization is the emergence of the
\textbf{centroid vector} $\bar{v}$. Each article $a_i$ assigned to
event $E$ contributes a 6-dimensional \texttt{strategic\_interest\_vector}
$v_i \in [-1, 1]^6$, representing its narrative perspective across the
dimensions of security, economy, technology, resource, ideology, and
environment. The centroid is defined as the arithmetic mean across all
contributing perspectives:

\begin{equation}
  \bar{v} = \frac{1}{n} \sum_{i=1}^{n} v_i
\end{equation}

The centroid represents the \textbf{gravitational center of the
event} — the best available approximation of the event's objective
geopolitical weight, computed from the ensemble of all contributing
narrative perspectives. As more articles converge on the same
\texttt{event\_id}, the centroid stabilizes: the variance of
$\{v_i\}$ decreases and the estimate of $\bar{v}$ becomes
progressively more reliable.

This gives rise to a natural measure of \textbf{crystallization
completeness}. When all six canonical narrative perspectives
(jp, us, cn, gb, eu, qa) have been compiled for a given event,
the \texttt{source\_count} reaches 6. At this point the event is
considered \emph{fully crystallized} — its centroid has been
computed from the maximum available set of viewpoints and no
further perspective shift is expected. The distance between any
individual narrative vector $v_i$ and the centroid $\bar{v}$,

\begin{equation}
  \text{CDC}_i = \| v_i - \bar{v} \|_2
\end{equation}

then functions as the \textbf{Cross-Domain Correlation (CDC)} metric:
a quantitative measure of how far a given national narrative deviates
from the ensemble center of gravity. High CDC values indicate
narrative outliers — perspectives that are structurally distant from
the consensus — and serve as the primary signal for the
\texttt{audit\_aura} indicator in the presentation layer.

\subsection{Narrative Compression}

The final stage of quantification is \textbf{Narrative Compression},
the process by which structured narrative data is refined into a
high-density cognitive representation. This operation can be understood
as an \textbf{active entropy reduction} within the information space.

Raw, real-world data is inherently high-entropy---it is noisy,
continuous, and lacks explicit meaning. Narrative compression provides
several key advantages:

\begin{itemize}[noitemsep]
  \item \textbf{Computational Efficiency:} Reasoning over discrete, compressed event units is far more efficient than processing raw linguistic tokens.
  \item \textbf{Long-Term Coherence:} Stripping away noise reveals the narrative arc, enabling tracking of long-term dynamics that LLMs often lose.
  \item \textbf{Storage and Retrieval:} Compression enables construction of vast, searchable Narrative Databases.
\end{itemize}

Ultimately, narrative compression transforms ``text'' into
``knowledge.'' This compressed, structured format serves as the
fundamental input for the Large Cognitive Model (LCM).

% ── 4. The Narrative Compiler ─────────────────────────────────────────────────
\section{The Narrative Compiler: Search and Optimization}

The Narrative Compiler does not merely extract data; it explores a vast
space of potential interpretations to find the most probable narrative
structures. Given that a single text can yield an exponential number of
event connections, the Compiler employs specific \textbf{Search
Constraints} to manage this complexity and generate the final
probability distribution $P(N \mid \text{text})$.

\subsection{Generating Possible Narrative Graphs}

The Compiler treats narrative construction as a graph-search problem.
From a set of extracted events, it attempts to build a \textbf{Narrative
Graph} $G$ where nodes represent events and edges represent causal or
temporal transitions.

\subsection{Beam Search for Hypothesis Management}

To maintain a manageable ``hash set'' of narrative states, the Compiler
utilizes \textbf{Beam Search}. Rather than exploring every possible
graph configuration, the Compiler keeps only the top-$k$ most plausible
narrative hypotheses at each step of the construction.

\begin{itemize}[noitemsep]
  \item \textbf{Top-$k$ Survival:} At each stage of event linking, only the graphs with the highest cumulative probability scores are retained.
  \item \textbf{Diversity Maintenance:} The beam search is tuned to preserve divergent narrative paths if their probabilities are sufficiently high.
\end{itemize}

\subsection{Graph Pruning and Information Entropy}

\textbf{Graph Pruning} is the active removal of low-probability or
redundant branches in the narrative tree.

\begin{itemize}[noitemsep]
  \item \textbf{Redundancy Filtering:} Multiple paths leading to the same structural state with lower efficiency are pruned.
  \item \textbf{Entropy Thresholding:} Branches that increase conditional entropy without adding semantic value are eliminated.
\end{itemize}

\subsection{Causal Constraints: The Logical Filter}

The most critical constraint is the \textbf{Causal Constraint}, which
acts as the ``physics'' of the Narrative Space. The Compiler applies
logical rules to validate every proposed edge in the graph:

\begin{itemize}[noitemsep]
  \item \textbf{Temporal Precedence:} A cause must strictly precede its effect.
  \item \textbf{Actor Consistency:} Actions must align with the capabilities and known motivations of the involved actors.
  \item \textbf{Interaction Logic:} The Action-Reaction loop must maintain coherent flow.
\end{itemize}

\subsection{Final State Synthesis}

Through these constraints, the Narrative Compiler collapses the
near-infinite possibilities of prose into a refined \textbf{hash set}.
This set represents the ``State'' of the narrative---a weighted
distribution of graphs that have survived beam search and passed all
causal filters.

% ── 5. The 24TAG Taxonomy ─────────────────────────────────────────────────────
\section{The 24TAG Taxonomy}

To ensure universal compatibility and deterministic hashing, the
Narrative Quantification protocol defines a fixed schema of 24
structured tags. These tags are categorized into four primary layers.

\begin{table}[h!]
\centering
\caption{The 24TAG Framework Definitions}
\small
\begin{tabular}{llll}
\toprule
\textbf{Layer} & \textbf{Tag ID} & \textbf{Tag Name} & \textbf{Description} \\
\midrule
\multirow{6}{*}{\textbf{Actor \& Agency}}
  & T01 & Primary\_Actor      & Main entity initiating the interaction \\
  & T02 & Secondary\_Actor    & Entity receiving or responding to action \\
  & T03 & Actor\_Role\_A      & Functional role of primary actor \\
  & T04 & Actor\_Role\_B      & Functional role of secondary actor \\
  & T05 & Actor\_Motivation\_A & Underlying goal for Actor A \\
  & T06 & Actor\_Motivation\_B & Underlying goal for Actor B \\
\midrule
\multirow{6}{*}{\textbf{Event \& Action}}
  & T07 & Action\_Type        & Category of event (Verbal, Physical, etc.) \\
  & T08 & Action\_Intensity   & Magnitude of action (0.0--1.0) \\
  & T09 & Target\_Resource    & Specific object or value being contested \\
  & T10 & Event\_Modality     & Factual, hypothetical, or desired \\
  & T11 & Temporal\_Sequence  & Relative order within narrative arc \\
  & T12 & Spatial\_Context    & Domain or environment of event \\
\midrule
\multirow{6}{*}{\textbf{Causal \& Logic}}
  & T13 & Causal\_Link\_Type  & Nature of connection (Direct Cause, etc.) \\
  & T14 & Conflict\_Nature    & Type of opposition (Internal, etc.) \\
  & T15 & Conflict\_Intensity & Severity of friction between actors \\
  & T16 & Resolution\_Status  & Degree of conflict settlement \\
  & T17 & Outcome\_Valence    & Positive or negative impact \\
  & T18 & Logic\_Consistency  & Alignment with prior narrative states \\
\midrule
\multirow{6}{*}{\textbf{Context \& Tone}}
  & T19 & Emotional\_Tone\_A  & Emotional state of Actor A \\
  & T20 & Emotional\_Tone\_B  & Emotional state of Actor B \\
  & T21 & Perspective\_Bias   & Point of view of the narrative \\
  & T22 & Info\_Asymmetry     & Knowledge gap between actors \\
  & T23 & Significance\_Score & Relative importance of event \\
  & T24 & Narrative\_Closure  & Extent of thread conclusion \\
\bottomrule
\end{tabular}
\end{table}

\subsection{Data Structure and Compilation}

Each tag ($T_{01 \dots 24}$) is populated by the Narrative Compiler
using the Beam Search and Causal Constraints defined in Section~4. The
resulting data structure is a JSON object. As specified in the
\textbf{Crystallization Protocol}, the \texttt{state\_hash} is
calculated by applying SHA-256 to the JCS-normalized string of these 24
tags:

\begin{lstlisting}[style=json]
{
  "T01": "Actor_A",
  "T02": "Actor_B",
  ...
  "T24": "Closure_Value",
  "state_hash": "6b86b273ff34fce19d6b804eff5a3f57..."
}
\end{lstlisting}

\begin{itemize}[noitemsep]
  \item \textbf{Normalization (JCS):} Following RFC~8785~\cite{rfc8785}, the Compiler sorts the 24 structured tags into a canonical JSON format.
  \item \textbf{Self-Referential Exclusion:} The \texttt{state\_hash} field itself is strictly excluded from the hash input to prevent circular dependencies.
  \item \textbf{Hashing (SHA-256):} The Compiler applies SHA-256 to the canonical byte-stream of the 24 tags.
\end{itemize}

The mathematical definition of the state is:

\begin{equation}
  state\_hash = \mathrm{SHA256}\!\left(\mathrm{JCS}\!\left(\{TAG_{01 \dots 24}\} \setminus \{state\_hash\}\right)\right)
\end{equation}

\subsection{Integrity and Verification}

By excluding the \texttt{state\_hash} from its own input, we create a
\textbf{One-Way Cognitive Seal}. A receiver can verify the integrity of
a narrative state by re-running the JCS and SHA-256 process on the
remaining 24 tags and comparing against the original \texttt{state\_hash}.

% ── 6. Narrative Space ────────────────────────────────────────────────────────
\section{Narrative Observation and the Narrative Space}

\subsection{Defining Narrative Space}

\begin{definition}
\textbf{Narrative Space} is the multi-dimensional space formed by the
set of possible narrative states. Each specific state is represented as
a point within this space.
\end{definition}

\subsection{The Mathematical Foundation: Graph Embedding}

Each narrative graph $G$ is projected into a $d$-dimensional
real-valued space via an embedding function $e$:

\begin{equation}
  e : G \rightarrow \mathbb{R}^d
\end{equation}

\subsection{Quantifying Divergence through Distance}

The divergence between two narrative states $G_1$ and $G_2$ is defined
by the distance between their corresponding embeddings:

\begin{equation}
  \text{distance} = \| e(G_1) - e(G_2) \|
\end{equation}

\begin{itemize}[noitemsep]
  \item \textbf{Low Distance (Convergence):} Indicates cognitive consensus.
  \item \textbf{High Distance (Divergence):} Indicates competing interpretations or ``Narrative Rifts.''
\end{itemize}

\subsection{Mapping the Narrative Ecosystem}

By observing the distribution and trajectory of vectors $e(G)$ within
the Narrative Space, researchers can identify:

\begin{itemize}[noitemsep]
  \item \textbf{Attractors:} Regions where narrative vectors cluster, indicating dominant societal beliefs.
  \item \textbf{Bifurcation Points:} Moments where a single event causes the distribution to split into two distant clusters.
  \item \textbf{Narrative Velocity:} The rate at which $e(G)$ moves through the space as new information is integrated.
\end{itemize}

% ── 7. Cognitive Infrastructure ───────────────────────────────────────────────
\section{Cognitive Infrastructure}

\subsection{From Knowledge Bases to Narrative Databases}

Traditional knowledge bases (e.g., Wikidata, Knowledge Graphs) are
effective at storing \emph{what} happened---discrete facts such as
dates, locations, and actors. However, they lack the capacity to store
\emph{how} those facts are interpreted. A \textbf{Narrative Database},
powered by the Narrative Compiler, stores \textbf{Crystallized Narrative
States}. This shifts the focus from ``Data Storage'' to ``Cognitive
State Storage.''

\subsection{Information Integrity and the State Hash}

The deterministic nature of the \texttt{state\_hash} provides a
mechanism for \textbf{Information Integrity}:

\begin{itemize}[noitemsep]
  \item \textbf{Proof of Interpretation:} A \texttt{state\_hash} attached to a report provides a ``Cognitive Seal.''
  \item \textbf{Divergence Detection:} Comparing \texttt{state\_hash} outputs across media outlets objectively indicates narrative structural differences.
\end{itemize}

\subsection{Applications of the Infrastructure}

\begin{itemize}[noitemsep]
  \item \textbf{Media \& Information Integrity:} Automated detection of narrative manipulation.
  \item \textbf{Geopolitical Narrative Mapping:} Real-time observation of how different nations interpret the same international event.
  \item \textbf{Corporate \& Economic Intelligence:} Tracking the evolution of market narratives through structured event formation.
\end{itemize}

% ── 8. LCM ────────────────────────────────────────────────────────────────────
\section{Reasoning over Narrative Distributions (LCM)}

By defining a state as $P(N \mid \text{text})$, the \textbf{Large
Cognitive Model (LCM)} gains the ability to perform ``superpositional''
reasoning.

\subsection{Inference as Distribution Transformation}

While LLMs predict tokens, the LCM performs inference by transforming
one probability distribution into another. The reasoning core evaluates
how the current distribution $P(N_t \mid \text{text}_t)$ evolves into a
future distribution $P(N_{t+1})$ when a new event occurs.

\subsection{Handling Ambiguity and Divergence}

This probabilistic approach allows the LCM to handle \textbf{Narrative
Divergence} natively. When a text is highly ambiguous, the Compiler
produces a broad ``hash set'' with distributed probabilities. The LCM
can track multiple ``possible worlds'' simultaneously.

\subsection{Reasoning over Causal Trajectories}

Because an event is defined as the \emph{minimum causal interaction
unit}, the LCM can track long-term narrative arcs without the
``forgetting'' or ``hallucination'' common in LLMs. Reasoning in an LCM
is not a search for the next word, but a calculation of \textbf{state
transitions}:

\begin{equation}
  \text{State } A\,(\mathrm{Hash}_1)
  \xrightarrow{\text{Causal Logic}}
  \text{State } B\,(\mathrm{Hash}_2)
\end{equation}

% ── 9. Human Memory ───────────────────────────────────────────────────────────
\section{Human Memory and Narrative Cognition}

Human memory does not function as a raw sensory log; rather, it
operates through a sophisticated process of \textbf{narrative
compression}.

\subsection{Memory as a Compiled Distribution}

When humans experience an event, the brain acts as a biological
\textbf{Narrative Compiler}, transforming continuous experience into a
discrete set of structured narrative units. Human memory can be viewed
as the storage of the probability distribution:

\begin{equation}
  P(N \mid \text{Experience})
\end{equation}

What we ``remember'' is not the raw text of reality, but the
\textbf{hash set} of the most probable narrative structures.

\subsection{The Primitives of Biological Storage}

The human cognitive system captures only essential semantic primitives:

\begin{itemize}[noitemsep]
  \item \textbf{Key Actors:} Identifying who was involved.
  \item \textbf{Causal Interpretations:} Assigning a ``why'' to the sequence of events.
  \item \textbf{Emotional Evaluation:} Tagging the event with a valence.
  \item \textbf{Outcomes:} The final state after the interaction is resolved.
\end{itemize}

\subsection{Biomimetic Validation}

The approach of Narrative Quantification is fundamentally
\textbf{biomimetic}. By using the \texttt{state\_hash} to represent a
``crystallized'' interpretation, we emulate the way human synapses
consolidate complex experiences into stable, retrievable ``concepts.''

% ── 10. Limitations ───────────────────────────────────────────────────────────
\section{Limitations}

\subsection{Subjectivity and Algorithmic Bias}

The primary challenge is the inherent subjectivity of human
interpretation. The ``priors'' used to assign probabilities are often
derived from the training data of the underlying models, which can
introduce systemic biases where the compiler favors specific cultural
tropes.

\subsection{Computational Complexity of Graph Generation}

Constructing \textbf{possible narrative graphs} is a non-trivial
computational task. In complex, multi-actor texts, the number of
potential causal connections grows exponentially.

\subsection{Cultural and Linguistic Heterogeneity}

Narrative structures are not universal. Different cultures employ
different storytelling archetypes. A 24TAG framework optimized for one
may fail to capture the nuances of another.

\subsection{Data Sovereignty and ``Narrative Warfare''}

If a centralized authority can define the ``standard'' \texttt{state\_hash}
for a global event, they possess the power to delegitimize divergent
but valid interpretations. The risk of using this infrastructure for
``narrative warfare'' is a serious ethical concern.

\subsection{Semantic Loss during Compression}

Narrative compression is an active entropy reduction process. There is a
persistent risk that subtle but crucial ``soft'' data---such as irony,
subtext, or poetic ambiguity---may be lost during the transition from
raw text to a structured \texttt{state\_hash}.

% ── 11. Conclusion ────────────────────────────────────────────────────────────
\section{Conclusion}

Despite these limitations, Narrative Quantification represents a
necessary evolution in Artificial Intelligence. By moving beyond the
probabilistic prediction of language tokens and toward the
\textbf{structured reasoning of narrative states}, we provide a
framework that more closely mirrors human cognition.

The transition from \textbf{Narrative Structures to Large Cognitive
Models (LCM)} offers more than just technical efficiency; it provides a
way to map the ``Narrative Space'' of our society, making the invisible
structures of our interpretations visible, measurable, and verifiable.

Future research will focus on mitigating bias through diverse
``compiler ensembles'' and exploring the decentralized governance of
narrative hashes to ensure that this cognitive infrastructure serves as
a tool for clarity, not control.

The path toward \textbf{Artificial General Intelligence (AGI)} lies not
in bigger models, but in deeper structures. Narrative Quantification is
the first step in building an AI that doesn't just talk like a human,
but understands like one.

% ── References ────────────────────────────────────────────────────────────────
\begin{thebibliography}{99}

\bibitem{rfc8785}
Rundgren, A., Jordan, B., \& Erdtman, S. (2020).
\textit{JSON Canonicalization Scheme (JCS)}.
RFC 8785. Internet Engineering Task Force.
\url{https://www.rfc-editor.org/rfc/rfc8785}

\bibitem{shannon1948}
Shannon, C. E. (1948).
A mathematical theory of communication.
\textit{Bell System Technical Journal}, 27(3), 379--423.

\bibitem{minsky1975}
Minsky, M. (1975).
A framework for representing knowledge.
In P. H. Winston (Ed.), \textit{The Psychology of Computer Vision} (pp. 211--277). McGraw-Hill.

\bibitem{schank1977}
Schank, R. C., \& Abelson, R. P. (1977).
\textit{Scripts, Plans, Goals and Understanding}.
Lawrence Erlbaum Associates.

\bibitem{bruner1991}
Bruner, J. (1991).
The narrative construction of reality.
\textit{Critical Inquiry}, 18(1), 1--21.

\bibitem{vaswani2017}
Vaswani, A., Shazeer, N., Parmar, N., et al. (2017).
Attention is all you need.
\textit{Advances in Neural Information Processing Systems}, 30.

\bibitem{devlin2018}
Devlin, J., Chang, M. W., Lee, K., \& Toutanova, K. (2018).
BERT: Pre-training of deep bidirectional transformers for language understanding.
\textit{arXiv preprint arXiv:1810.04805}.

\bibitem{hamilton2017}
Hamilton, W., Ying, Z., \& Leskovec, J. (2017).
Inductive representation learning on large graphs.
\textit{Advances in Neural Information Processing Systems}, 30.

\bibitem{ecma262}
Ecma International. (2024).
\textit{ECMAScript 2024 Language Specification (ECMA-262)}.
\url{https://tc39.es/ecma262/}

\end{thebibliography}

\end{document}
